
\subsection{Метеоданные}

\epigraph{
    You don't need a weatherman \\
    to know which way the wind blows
}{\textit{-- Bob Dylan}\\"Subterranean Homesick Blues" (1965)}


С раннего детства мы наблюдаем за окружающим нас миром.
Одна из граней этого процесса -- наблюдения за изменяющейся вокруг нас погодой.
Мы легко ощущаем это в моменте, но сколько труда стоит сохранить всю эту информацию и сделать ее общедоступной?
Насколько глубоко сохраняется эта память?
Как человечество рефлексировало ее?
Рассмотрим этот вопрос с исторической стороны!

\subsubsection{История}

Трепет и страх человека перед природой сопровождал его на протяжении всей истории.
Катастрофические проявления погодных феноменов сформировали по всему миру мономифы о великих потопах и бесконечных зимах.
Тем не менее, первые реальные письменные источники, содержащие описания погодных явлений дошли до нас из Греции V-IV веков до нашей эры.
Трактат Аристотеля «Метеорологика» заложил теоретические основы метеорологии, однако развитие ее как науки произошло лишь в XVII веке с рождением экспериментальной физики.
В 1654 году Фердинандо II Медичи под эгидой Флорентийской академии организовал первую в истории метеорологическую сеть из 10 станций размещенных от Флоренции до Варшавы и ставшую прообразом всех будущих наблюдательных систем.
Термоскоп Галилея, барометр Торричелли и труды Паскаля позволили перейти от качественных описаний к количественным измерениям, заменив бытовые «приметы» обощениями метематического анализа.

В России систематические наблюдения начали выполняться по указу Петра I с 1722 года в Санкт-Петербурге. 
В 1733 году экспедиция Беринга создала цепочку метеостанций от Казани до Сибири. 
В 1780 году по инициативе курфюрста Карла Теодора была создана "Societas Meteorologica Palatina" -- первая в мире международная метеорологическая сеть в современном смысле, объединившая 39 стандартно оснащенных станций от Европы до Северной Америки.

Все эти силы были согласованы целью обеспечении безопасности морского судоходства.
Черноморская буря, разметавшая и уничтожившая во время Крымской войны французский и британский флоты, стала достаточным обоснованием к созданию первой синоптической карты погоды, обновляемой в реальном времени, ставшей возможной благодаря появлению телеграфной связи.
В 1873 году состоялся Первый международный метеорологический конгресс, установивший единые сроки измерений и стандарты передачи и обмена данными.
Все это стало прообразом современных глобальных климатических моделей, данные наследников которых используются в представляемой работе.

На рубеже XIX–XX веков норвежский геофизик Вильгельм Бьеркнес  сформулировал тезис, что точный метеопрогноз возможен лишь при анализе глобальной циркуляции атмосферы.
Вслед за этим в 1922 году британский математик Льюис Фрай Ричардсон опубликовал книгу "Weather Prediction by Numerical Process", в которой предложил вычислять прогноз погоды путем численного интегрирования уравнений гидродинамики атмосферы. 

Ричардсон мечтал о «фабрике прогнозов» -- зале из 64 тысяч счетчиков, работающих параллельно.
Его методология оказалась верной, но опередила свое время и не смогла быть реализована в полной мере~\cite{Kimura2002NumericalWeatherPrediction}.
В 1950 году на компьютере ENIAC в США метеорологами Джулом Чарни, Рагнаром Фьортофтом и математиком Джоном фон Нейманом был впервые получен численный прогноз погоды с помощью ЭВМ. 
Расчет суточного прогноза занял 24 часа вычислительного времени, но это был поворотный момент в истории метеорологии. 

К середине XX века стало очевидно, что точный прогноз погоды требует не разрозненных национальных сетей, а глобальной системы наблюдений и обмена данными. 
Атмосферные процессы не знают границ: циклон, зародившийся над Атлантикой, через трое суток достигает Европы, а для его прогноза нужны данные со всего океана. 
23 марта 1950 года вступила в силу Конвенция Всемирной Метеорологической Организации (ВМО) ставшей в декабре 1951 года ВМО стала специализированным агентством ООН.

Создание ВМО заложило институциональную основу, но технологический прорыв обеспечили два события: запуск первого метеоспутника TIROS-1 в 1960 году и развертывание программы World Weather Watch с 1963 года.
1 апреля 1960 года спутник TIROS-1 передал первые 14 000 снимков облачности из космоса, впервые показав синоптикам погоду над океанами и полюсами. 
Это стало катализатором создания трехуровневой мировой системы метеоцентров с единой телекоммуникацией, связавшей более 10 тысяч наземных станций и десятки спутников. 
Параллельно суперкомпьютеры превратили мечту Ричардсона в реальность.
Уже в 1979 году ECMWF выдавал 10-дневные глобальные прогнозы, а сегодня ИИ-модели воспроизводят динамику всей планетарной атмосферы за минуты.

В 1964 году в рамках Всемирной службы погоды в СССР был создан Мировой метеорологический центр, в 1965-м году он объединился с Центральным институтом прогнозов в Гидрометцентр СССР с функциями мирового и регионального центров ВМО. 
После распада СССР в 1992 году на базе советского Госкомгидромета был создан Комитет по гидрометеорологии и мониторингу окружающей среды России.

Сегодня в России основным источником метеорологических данных является Федеральная служба по гидрометеорологии и мониторингу окружающей среды (Росгидромет). 
Правовая основа деятельности службы закреплена в Федеральном законе № 113-ФЗ от 19.07.1998 «О гидрометеорологической службе», который устанавливает принцип открытости и общедоступности информации о состоянии окружающей среды (статья 14)~\cite{rf_fedlaw_113fz_1998}. Приказ Минприроды России № 387 от 11.07.2025~\cite{minprirody_order_387_2025}, вступивший в силу 1 марта 2026 года, конкретизирует состав информации общего назначения и порядок ее предоставления.

Несмотря на декларируемую доступность~\cite{minprirody_order_387_2025}, механизм получения данных предполагает направление официального запроса с рассмотрением в срок до 30 дней.
Такой порядок, постулируя открытость де-юре, де-факто создает административные барьеры для оперативного использования данных в исследовательских и практических целях. 
Это определяет необходимость обращения к альтернативным источникам, предоставляющим свободный доступным через открытый API без согласований.

Таким образом, для научных исследований прибрежной гидродинамики и волновой экспозиции практически более эффективными оказываются международные открытые платформы, нежели национальные бюрократические процедуры запроса данных у Росгидромета.
